% !TeX program = xelatex
\documentclass[UTF8,11pt,a4paper,linespread=1]{ctexart}
\usepackage[margin=25mm]{geometry}
\usepackage{amsmath,amssymb,bm}
\usepackage{setspace}
\setstretch{1.35}

\begin{document}
\pagestyle{plain}
\pagenumbering{arabic}
\begin{center}
  {\zihao{3}\bfseries 面向校园共享单车的需求预测与调度优化}\par
  \vspace{.7em}
  {\large\bfseries 摘\quad 要}
\end{center}

针对校园共享单车在高峰时段供需不均的问题，本文以站点借还记录与时段特征为输入，
依次处理需求预测、调度方案生成和方案检验三个任务。

对于需求预测，采用［填写实际模型名称］刻画站点与时段的关系，
以［填写验证方式］划分训练数据与测试数据，并使用平均绝对误差衡量预测效果。
测试结果为［填写实际数值及单位］，与［填写基准方法］相比的变化为［填写实际数值］。

对于调度优化，以［填写实际目标函数］为目标，考虑车辆容量、站点库存与调度时限约束，
采用［填写实际求解方法］得到可行方案。方案在［填写场景］下的结果为［填写实际指标］。

通过［填写实际进行的敏感性或稳健性检验］检查方案对关键参数的依赖。
结果表明［填写由实验支持的结论］。模型尚未纳入［填写具体限制］，
应用到其他校园时，需要重新估计相关参数。

\par\medskip\noindent
\textbf{关键词：}需求预测；调度优化；约束规划；模型检验

\end{document}
